% !TeX program = pdflatex
\documentclass[10pt, aspectratio=169]{beamer}

% Preámbulo modular para presentaciones Beamer (Modelación Estadística)
% Diseño y paleta institucional del Tecnológico de Monterrey

\documentclass[aspectratio=169,xcolor={dvipsnames,table}]{beamer}

\usepackage[utf8]{inputenc}
\usepackage[spanish,mexico]{babel}
\usepackage{amsmath,amssymb,amsfonts,mathtools}
\usepackage{graphicx}
\usepackage{listings}
\usepackage{booktabs}
\usepackage{tikz}
\usetikzlibrary{matrix,arrows,decorations.pathmorphing,shapes.geometric,positioning}

% Paleta Institucional Tecnológico de Monterrey
\definecolor{TecAzul}{HTML}{0039A6}       % Azul TEC: color primario institucional
\definecolor{TecAzulOscuro}{HTML}{1A2E51} % Azul marino oscuro
\definecolor{TecRojo}{HTML}{EC2661}       % Rojo de acento (paleta miTec)
\definecolor{TecGris}{HTML}{646464}       % Gris neutro para comentarios y subtítulos
\definecolor{TecFondoBloque}{HTML}{F4F6F9}% Fondo suave para bloques

% Configuración de Tema Beamer
\usetheme{Boadilla}
\usecolortheme[named=TecAzulOscuro]{structure}
\setbeamercolor{alerted text}{fg=TecRojo}
\setbeamercolor{palette primary}{bg=TecAzulOscuro,fg=white}
\setbeamercolor{palette secondary}{bg=TecAzul,fg=white}
\setbeamercolor{palette tertiary}{bg=TecAzulOscuro!80!black,fg=white}
\setbeamercolor{title}{fg=TecAzulOscuro,bg=TecFondoBloque}
\setbeamercolor{frametitle}{fg=TecAzulOscuro,bg=TecFondoBloque}

% Bloques personalizados elegantes
\setbeamercolor{block title}{bg=TecAzulOscuro,fg=white}
\setbeamercolor{block body}{bg=TecFondoBloque,fg=black}
\setbeamercolor{block title alerted}{bg=TecRojo,fg=white}
\setbeamercolor{block body alerted}{bg=TecRojo!8,fg=black}
\setbeamercolor{block title example}{bg=TecAzul,fg=white}
\setbeamercolor{block body example}{bg=TecAzul!8,fg=black}

% Redondear bordes y añadir sombra sutil
\setbeamertemplate{blocks}[rounded][shadow=true]
\setbeamertemplate{navigation symbols}{}

% Pie de página limpio con número de diapositiva
\setbeamertemplate{footline}{
  \leavevmode%
  \hbox{%
  \begin{beamercolorbox}[wd=.7\paperwidth,ht=2.25ex,dp=1ex,left]{author in head/foot}%
    \hspace*{2ex}\insertshortauthor~~(\insertshortinstitute) \hfill \insertshorttitle
  \end{beamercolorbox}%
  \begin{beamercolorbox}[wd=.3\paperwidth,ht=2.25ex,dp=1ex,right]{date in head/foot}%
    \insertshortdate{}\hspace*{2em}
    \insertframenumber{} / \inserttotalframenumber\hspace*{2ex}
  \end{beamercolorbox}}%
  \vskip0pt%
}

% Configuración de Listings de Python para visualización pedagógica de código
\lstset{
    language=Python,
    basicstyle=\ttfamily\footnotesize,
    keywordstyle=\color{TecAzulOscuro}\bfseries,
    stringstyle=\color{TecRojo},
    commentstyle=\color{TecGris}\itshape,
    showstringspaces=false,
    breaklines=true,
    frame=lines,
    rulecolor=\color{TecAzulOscuro!30},
    backgroundcolor=\color{TecFondoBloque},
    aboveskip=0.5em,
    belowskip=0.5em,
    literate={á}{{\'a}}1 {é}{{\'e}}1 {í}{{\'i}}1 {ó}{{\'o}}1 {ú}{{\'u}}1
             {Á}{{\'A}}1 {É}{{\'E}}1 {Í}{{\'I}}1 {Ó}{{\'O}}1 {Ú}{{\'U}}1
             {ñ}{{\~n}}1 {Ñ}{{\~N}}1 {¿}{{?`}}1 {¡}{{!`}}1
}

% Metadatos institucionales por defecto
\title[Modelación Estadística]{Modelación Estadística para Ciencia de Datos e Ingeniería}
\author[J. Castillo Colmenares]{Juliho David Castillo Colmenares}
\institute[Tec de Monterrey]{Tecnológico de Monterrey \\ Escuela de Ingeniería y Ciencias}
\date{\today}

% Comandos y macros estadísticas compartidas para las presentaciones Beamer (Metropolis)

% Conjuntos numéricos básicos
\newcommand{\R}{\mathbb{R}}
\newcommand{\N}{\mathbb{N}}
\newcommand{\Q}{\mathbb{Q}}
\newcommand{\Z}{\mathbb{Z}}
\newcommand{\set}[1]{\left\{ #1 \right\}}
\newcommand{\abs}[1]{\left|#1\right|}
\newcommand{\norm}[1]{\left\| #1 \right\|}

% Comandos estadísticos y probabilísticos del curso
\newcommand{\Var}{\operatorname{Var}}
\newcommand{\cov}{\operatorname{Cov}}
\newcommand{\corr}{\rho}
\newcommand{\comb}[2]{\begin{pmatrix} #1 \\ #2 \end{pmatrix}}
\newcommand{\s}{\sigma}
\newcommand{\particion}{\mathfrak{P}}
\newcommand{\card}[1]{\#\left( #1 \right)}
\newcommand{\seq}[3]{\set{#1}_{#2}^{#3}}
\newcommand{\E}{\mathbb{E}}
\newcommand{\Prob}{\mathbb{P}}

% Entornos pedagógicos y cajas en español académico natural
\newenvironment{discusion}{%
  \begin{alertblock}{Pregunta de Discusión}%
}{%
  \end{alertblock}%
}

\newenvironment{reflexion}{%
  \begin{alertblock}{Reflexión para el Aula}%
}{%
  \end{alertblock}%
}

\newenvironment{paradoja}{%
  \begin{alertblock}{Paradoja de Exploración}%
}{%
  \end{alertblock}%
}

\newenvironment{motivacion}{%
  \begin{exampleblock}{Motivación: Ciencia de Datos en la Práctica}%
}{%
  \end{exampleblock}%
}

\newenvironment{retoclase}{%
  \begin{block}{Reto Rápido en Equipo}%
}{%
  \end{block}%
}


\title[Teorema de Bayes]{Sección 02.05: Teorema de Bayes e Inversión Condicional}
\subtitle{De la Verosimilitud a la Probabilidad a Posteriori en Ingeniería y Diagnóstico}
\author[J. Castillo Colmenares]{Juliho Castillo Colmenares}
\institute{Tecnológico de Monterrey}
\date{\vspace{-2.0cm}}

\begin{document}

% Slide 01: Portada
\begin{frame}[plain]
  \vspace{-0.5cm}
  \titlepage
\end{frame}

% Slide 02: Hoja de Ruta
\begin{frame}{Hoja de Ruta --- Capítulo 02: Teoría de la Probabilidad}
  \begin{block}{Estructura Curricular de la Unidad 1}
    \footnotesize
    \begin{enumerate}\itemsep=0.08cm
      \item \textcolor{gray}{Sección 02.01: Introducción a la Probabilidad y Fenómenos Aleatorios}
      \item \textcolor{gray}{Sección 02.02: Teoría de Conjuntos y Axiomas de Kolmogorov}
      \item \textcolor{gray}{Sección 02.03: Fundamentos de Probabilidad y Técnicas de Conteo}
      \item \textcolor{gray}{Sección 02.04: Probabilidad Condicional e Independencia Estadística}
      \item \textbf{\textcolor{TecRojo}{Sección 02.05: Teorema de Bayes e Inversión Condicional}} $\leftarrow$ \emph{Foco actual}
      \item \textcolor{gray}{Sección 02.06: Muestreo Aleatorio y Teorema del Límite Central (TLC)}
    \end{enumerate}
  \end{block}
  \vspace{-0.05cm}
  \pause
  \begin{alertblock}{Objetivo Didáctico de la Sección 02.05}
    \footnotesize
    Dominar el álgebra de la inversión probabilística para transformar verosimilitudes empíricas en probabilidades a posteriori, analizar la falacia de la tasa base y modelar sistemas de actualización bayesiana secuencial en ingeniería y diagnóstico técnico.
  \end{alertblock}
\end{frame}

% Slide 03: Motivación — La Inversión de Condicionales
\begin{frame}{Motivación --- La Inversión de Condicionales}
  \begin{columns}[T]
    \begin{column}{0.48\textwidth}
      \begin{block}{Dirección Directa (Verosimilitud)}
        \small
        \begin{itemize}\itemsep=0.1cm
          \item En modelación física e industrial, conocemos la probabilidad de observar un \textbf{efecto o evidencia} $B$ dada una causa o hipótesis subyacente $A$:
          $$P(\text{Efecto} \mid \text{Causa}) = P(B \mid A)$$
          \item \emph{Ejemplo:} La probabilidad de que un sensor de vibración emite una alarma ($B$) cuando una turbina presenta falla mecánica ($A$).
        \end{itemize}
      \end{block}
    \end{column}
    \begin{column}{0.48\textwidth}
      \begin{alertblock}{Dirección Inversa (Posteriori)}
        \small
        \begin{itemize}\itemsep=0.1cm
          \item En diagnóstico médico, robótica y ciencia de datos, el problema real opera en sentido inverso: observamos la evidencia $B$ y debemos inferir la causa $A$:
          $$P(\text{Causa} \mid \text{Efecto}) = P(A \mid B)$$
          \item \emph{Ejemplo:} Si la alarma sonó ($B$), ¿cuál es la probabilidad real de una falla catastrófica ($A$)?
        \end{itemize}
      \end{alertblock}
    \end{column}
  \end{columns}
\end{frame}

% Slide 04: Deducción del Teorema de Bayes en Forma Simple
\begin{frame}{Deducción del Teorema de Bayes en Forma Simple}
  \fontsize{7.5pt}{9pt}\selectfont
  \begin{block}{Simetría de la Intersección}
    \fontsize{7.5pt}{9pt}\selectfont
    Sean $A$ y $B$ dos eventos en el espacio muestral $S$ con probabilidades estrictamente positivas ($P(A)>0, P(B)>0$). Por la regla de multiplicación de probabilidad condicional, la probabilidad conjunta $P(A \cap B)$ es algebraicamente simétrica:
    $$P(A \cap B) = P(B \mid A)P(A) = P(A \mid B)P(B)$$
    Despejando directamente $P(A \mid B)$ obtenemos la formulación elemental del \textbf{Teorema de Bayes}:
    $$P(A \mid B) = \frac{P(B \mid A)P(A)}{P(B)}$$
  \end{block}
  \vspace{-0.15cm}
  \pause
  \begin{alertblock}{Anatomía y Forma en Razón de Momios (Odds Form)}
    \fontsize{7.5pt}{9pt}\selectfont
    \begin{itemize}\itemsep=0.08cm
      \item $P(A)$: \textbf{Probabilidad a Priori} (creencia o conocimiento inicial antes de observar la evidencia).
      \item $P(B \mid A)$: \textbf{Verosimilitud} (\emph{Likelihood}, compatibilidad de la evidencia con la hipótesis).
      \item $P(A \mid B)$: \textbf{Probabilidad a Posteriori} (creencia actualizada tras incorporar el dato $B$).
      \item \textbf{Formulacion en Momios:} Para dos hipótesis rivales $A_1, A_2$: $\displaystyle \frac{P(A_1 \mid B)}{P(A_2 \mid B)} = \frac{P(B \mid A_1)}{P(B \mid A_2)} \cdot \frac{P(A_1)}{P(A_2)}$.
    \end{itemize}
  \end{alertblock}
\end{frame}

% Slide 05: Teorema de Bayes para Particiones Exhaustivas
\begin{frame}{Teorema de Bayes para Particiones Exhaustivas}
  \fontsize{7.5pt}{9pt}\selectfont
  \begin{block}{Expansión mediante el Teorema de Probabilidad Total}
    \fontsize{7.5pt}{9pt}\selectfont
    En problemas complejos de ingeniería, el espacio muestral está dividido en una \textbf{partición finita y exhaustiva} de causas mutuamente excluyentes $\{B_1, B_2, \ldots, B_k\}$ tal que $S = \bigcup_{i=1}^k B_i$. Para evaluar la probabilidad de cualquier estado particular $B_j$ dado el evento observable $A$, expandimos el denominador $P(A)$:
    $$P(B_j \mid A) = \frac{P(A \mid B_j)P(B_j)}{\sum_{i=1}^k P(A \mid B_i)P(B_i)}, \quad \text{para todo } j \in \{1, 2, \ldots, k\}$$
  \end{block}
  \vspace{-0.15cm}
  \pause
  \begin{alertblock}{Interpretación del Denominador (Constante de Normalización)}
    \fontsize{7.5pt}{9pt}\selectfont
    El denominador $\sum_{i=1}^k P(A \mid B_i)P(B_i)$ representa la \textbf{probabilidad marginal del evento observable} $P(A)$. Actúa como una constante de normalización rigurosa que garantiza que la suma de las probabilidades a posteriori de todas las causas posibles sea exactamente igual a la unidad: $\sum_{j=1}^k P(B_j \mid A) = 1$.
  \end{alertblock}
\end{frame}

% Slide 06: Diagnóstico Clínico y Falacia de la Tasa Base
\begin{frame}{Diagnóstico Clínico y la Falacia de la Tasa Base}
  \begin{columns}[T]
    \begin{column}{0.48\textwidth}
      \begin{block}{Métricas Condicionales del Test}
        \small
        Sea $E$ la presencia de una enfermedad y $T^+$ un resultado clínico positivo:
        \begin{itemize}\itemsep=0.1cm
          \item \textbf{Sensibilidad ($P(T^+ \mid E)$):} Probabilidad de detectar la enfermedad cuando está presente.
          \item \textbf{Especificidad ($P(T^- \mid E')$):} Probabilidad de dar negativo cuando el paciente está sano.
          \item \textbf{Tasa de Falsos Positivos:} $P(T^+ \mid E') = 1 - \text{Especificidad}$.
        \end{itemize}
      \end{block}
    \end{column}
    \begin{column}{0.48\textwidth}
      \begin{alertblock}{La Falacia de la Tasa Base (Prevalencia)}
        \small
        \begin{itemize}\itemsep=0.1cm
          \item El \textbf{Valor Predictivo Positivo (VPP)} es:
          $$P(E \mid T^+) = \frac{P(T^+ \mid E)P(E)}{P(T^+)}$$
          \item Si la prevalencia $P(E)$ es extremadamente baja ($1\%$), el volumen de falsos positivos en la población sana abruma a los verdaderos positivos, haciendo que $P(E \mid T^+) \ll 50\%$ incluso con alta sensibilidad.
        \end{itemize}
      \end{alertblock}
    \end{column}
  \end{columns}
\end{frame}

% Slide 07: Actualización Bayesiana Secuencial
\begin{frame}{Actualización Bayesiana Secuencial en Ingeniería}
  \fontsize{7.5pt}{9pt}\selectfont
  \begin{block}{El Posterior de Hoy es la Distribución a Priori de Mañana}
    \fontsize{7.5pt}{9pt}\selectfont
    Cuando un sistema industrial o robot autónomo recibe un flujo continuo de observaciones independientes dadas las hipótesis ($B_1, B_2, \ldots, B_t$), la inferencia bayesiana opera de forma recursiva sin necesidad de reevaluar todo el historial desde cero:
    $$P(A \mid B_1) = \frac{P(B_1 \mid A)P(A)}{P(B_1)} \implies P(A \mid B_1 \cap B_2) = \frac{P(B_2 \mid A)P(A \mid B_1)}{P(B_2 \mid B_1)}$$
  \end{block}
  \vspace{-0.15cm}
  \pause
  \begin{alertblock}{Convergencia Hacia la Verdad en Robótica e IA}
    \fontsize{7.5pt}{9pt}\selectfont
    \begin{itemize}\itemsep=0.08cm
      \item En cada iteración $t$, el nuevo dato $B_t$ ajusta el factor de verosimilitud para actualizar la creencia sobre el estado del sistema $\theta$.
      \item Esta arquitectura recursiva constituye el motor matemático de los \textbf{Filtros de Kalman}, la localización en robótica móvil (SLAM) y el aprendizaje de parámetros en modelos probabilísticos modernos.
    \end{itemize}
  \end{alertblock}
\end{frame}

% Slide 08: Clasificador Naive Bayes en Alta Dimensionalidad
\begin{frame}{Clasificador Naive Bayes en Alta Dimensionalidad}
  \fontsize{7.5pt}{9pt}\selectfont
  \begin{block}{El Supuesto de Independencia Condicional (Naive Bayes)}
    \fontsize{7.5pt}{9pt}\selectfont
    En clasificación de textos (ej. filtros anti-spam) o diagnóstico con múltiples síntomas $W = W_1 \cap W_2 \cap \ldots \cap W_m$, calcular la probabilidad conjunta exacta $P(W \mid C)$ es computacionalmente intratable. El algoritmo \textbf{Naive Bayes} asume que los atributos son condicionalmente independientes dada la clase $C$:
    $$P(W_1 \cap W_2 \cap \ldots \cap W_m \mid C) = \prod_{i=1}^m P(W_i \mid C)$$
  \end{block}
  \vspace{-0.15cm}
  \pause
  \begin{alertblock}{Regla de Decisión A Posteriori}
    \fontsize{7.5pt}{9pt}\selectfont
    La probabilidad de pertenencia a una clase $C_k$ dado el vector de atributos observados $W$ se calcula mediante el producto de verosimilitudes marginales:
    $$P(C_k \mid W) = \frac{P(C_k) \prod_{i=1}^m P(W_i \mid C_k)}{\sum_{j} P(C_j) \prod_{i=1}^m P(W_i \mid C_j)}$$
    Pese al supuesto simplificador de independencia, Naive Bayes exhibe una precisión extraordinaria en sistemas de decisión en tiempo real.
  \end{alertblock}
\end{frame}

% Slide 09A: Laboratorio Python — Test Médico e Inversión
\begin{frame}[fragile]{Laboratorio en Python --- Test Médico e Inversión (1/4)}
  \vspace{-0.1cm}
  \begin{block}{Código Fuente: \texttt{simulate\_medical\_test} (\texttt{02.05\_bayes\_theorem.py})}
    \lstinputlisting[language=Python, basicstyle=\fontsize{4.6pt}{5.4pt}\selectfont\ttfamily, firstline=16, lastline=43]{../../code/02_teoria_probabilidad/02.05_bayes_theorem.py}
  \end{block}
\end{frame}

% Slide 09B: Laboratorio Python — Filtro Naive Bayes
\begin{frame}[fragile]{Laboratorio en Python --- Filtro Naive Bayes (2/4)}
  \vspace{-0.1cm}
  \begin{block}{Código Fuente: \texttt{evaluate\_naive\_bayes\_spam} (\texttt{02.05\_bayes\_theorem.py})}
    \lstinputlisting[language=Python, basicstyle=\fontsize{5.5pt}{6.5pt}\selectfont\ttfamily, firstline=48, lastline=63]{../../code/02_teoria_probabilidad/02.05_bayes_theorem.py}
  \end{block}
\end{frame}

% Slide 09C: Laboratorio Python — Actualización Secuencial
\begin{frame}[fragile]{Laboratorio en Python --- Actualización Secuencial (3/4)}
  \vspace{-0.1cm}
  \begin{block}{Código Fuente: \texttt{sequential\_coin\_updating} (\texttt{02.05\_bayes\_theorem.py})}
    \lstinputlisting[language=Python, basicstyle=\fontsize{4.6pt}{5.4pt}\selectfont\ttfamily, firstline=68, lastline=93]{../../code/02_teoria_probabilidad/02.05_bayes_theorem.py}
  \end{block}
\end{frame}

% Slide 09D: Laboratorio Python — Salida en Consola Verificada
\begin{frame}[fragile]{Laboratorio en Python --- Salida en Consola Verificada (4/4)}
  \vspace{-0.1cm}
  \begin{columns}[T]
    \begin{column}{0.54\textwidth}
      \begin{block}{1. Test Médico \& 2. Naive Bayes}
        \fontsize{4.4pt}{5.2pt}\selectfont
\begin{verbatim}
--- 1. Medical Diagnostic Test (Base Rate Fallacy) ---
Total Test Positives:  11,667 / 100,000
True Positives:        1,891 (True Diseased among Positives)
Empirical PPV (Sim):   0.1621 (16.21%)
Theoretical PPV:       0.1624 (16.24%)

--- 2. Naive Bayes Spam Filter (Keywords W1, W2, W3) ---
Joint Likelihood P(W|Spam):  0.33600
Joint Likelihood P(W|Legit): 0.00200
Marginal Evidence P(W):      0.20240
Posterior P(Spam | W):       0.996047 (99.60%)
\end{verbatim}
      \end{block}
    \end{column}
    \begin{column}{0.44\textwidth}
      \begin{block}{3. Actualización Secuencial}
        \fontsize{4.4pt}{5.2pt}\selectfont
\begin{verbatim}
--- 3. Sequential Bayesian Updating ---
Coin Flips: [H, H, T, H, H]
Step 0 | P(Biased Coin = 0.80): 0.5000 (50.00%)
Step 1 | P(Biased Coin = 0.80): 0.6154 (61.54%)
Step 2 | P(Biased Coin = 0.80): 0.7191 (71.91%)
Step 3 | P(Biased Coin = 0.80): 0.5059 (50.59%)
Step 4 | P(Biased Coin = 0.80): 0.6210 (62.10%)
Step 5 | P(Biased Coin = 0.80): 0.7239 (72.39%)
\end{verbatim}
      \end{block}
    \end{column}
  \end{columns}
\end{frame}

% Slide 10A: Ejercicio en Clase — Nivel 1: Fundamental (Enunciado)
\begin{frame}{Ejercicio en Clase --- Nivel 1: Fundamental (1/2)}
  \begin{block}{Problema 2.5.4 --- Estudio y Probabilidad a Posteriori de Éxito}
    \footnotesize
    Un estudiante universitario se prepara para presentar una evaluación integradora general. La experiencia estadística indica que:
    \begin{itemize}\itemsep=0.1cm
      \item Si estudia debidamente ($E$), la probabilidad condicional de aprobar ($A$) es del $90\%$: $P(A \mid E) = 0.90$.
      \item Si no estudia ($E'$), la probabilidad condicional de aprobar es de apenas el $30\%$: $P(A \mid E') = 0.30$.
      \item La probabilidad a priori de que el estudiante estudie de manera rigurosa es del $70\%$: $P(E) = 0.70$.
    \end{itemize}
    \begin{enumerate}\itemsep=0.15cm
      \item Calcula la probabilidad total de que el estudiante apruebe el examen ($P(A)$).
      \item Si al publicarse las calificaciones se observa que el estudiante aprobó ($A$), ¿cuál es la probabilidad condicional de que haya estudiado ($P(E \mid A)$)?
    \end{enumerate}
  \end{block}
\end{frame}

% Slide 10B: Ejercicio en Clase — Nivel 1: Fundamental (Resolución)
\begin{frame}{Ejercicio en Clase --- Nivel 1: Fundamental (2/2)}
  \begin{columns}[T]
    \begin{column}{0.48\textwidth}
      \begin{block}{1. Probabilidad Total de Aprobar ($P(A)$)}
        \fontsize{6.5pt}{7.8pt}\selectfont
        Las hipótesis exhaustivas son $\{E, E'\}$ con $P(E)=0.70$ y $P(E')=0.30$. Por el Teorema de Probabilidad Total:
        $$P(A) = P(A \mid E)P(E) + P(A \mid E')P(E')$$
        Sustituyendo los valores del problema:
        $$P(A) = (0.90)(0.70) + (0.30)(0.30)$$
        $$P(A) = 0.63 + 0.09 = \mathbf{0.72 \text{ (72\%)}}$$
      \end{block}
    \end{column}
    \begin{column}{0.48\textwidth}
      \begin{alertblock}{2. Inversión Bayesiana ($P(E \mid A)$)}
        \fontsize{6.5pt}{7.8pt}\selectfont
        Aplicando el Teorema de Bayes para invertir la dirección condicional:
        $$P(E \mid A) = \frac{P(A \mid E)P(E)}{P(A)}$$
        Sustituyendo el numerador y el denominador ya calculado:
        $$P(E \mid A) = \frac{0.63}{0.72} = \frac{7}{8} = \mathbf{0.875 \text{ (87.5\%)}}$$
      \end{alertblock}
    \end{column}
  \end{columns}
  \vspace{-0.1cm}
  \pause
  \begin{exampleblock}{Interpretación Pedagógica}
    \fontsize{6.5pt}{7.8pt}\selectfont
    Sabiendo que el estudiante aprobó, nuestra certeza sobre su esfuerzo académico aumenta significativamente del $70\%$ a priori al $87.5\%$ a posteriori. El éxito observado constituye una evidencia sólida a favor de la preparación constante.
  \end{exampleblock}
\end{frame}

% Slide 11A: Ejercicio en Clase — Nivel 2: Operativo (Enunciado)
\begin{frame}{Ejercicio en Clase --- Nivel 2: Operativo (1/2)}
  \begin{block}{Problema 2.5.2 --- Evaluación de Test Médico Diagnóstico (VPP y VPN)}
    \footnotesize
    Un laboratorio clínico desarrolla una nueva prueba molecular para detectar una patología de baja prevalencia. Las características epidemiológicas y técnicas del test son las siguientes:
    \begin{itemize}\itemsep=0.1cm
      \item \textbf{Prevalencia a priori:} En la población general, el $2\%$ de las personas porta la enfermedad ($P(E) = 0.02$).
      \item \textbf{Sensibilidad:} Si el paciente está enfermo, la probabilidad de que el test dé positivo es de $0.95$ ($P(+ \mid E) = 0.95$).
      \item \textbf{Especificidad:} Si el paciente está sano ($E'$), la probabilidad de que el test dé negativo es de $0.90$ ($P(- \mid E') = 0.90$). Por lo tanto, la tasa de falsos positivos es $P(+ \mid E') = 0.10$.
    \end{itemize}
    \begin{enumerate}\itemsep=0.15cm
      \item Calcula el Valor Predictivo Positivo (VPP: $P(E \mid +)$) y explica por qué es sorprendentemente bajo.
      \item Calcula el Valor Predictivo Negativo (VPN: $P(E' \mid -)$).
    \end{enumerate}
  \end{block}
\end{frame}

% Slide 11B: Ejercicio en Clase — Nivel 2: Operativo (Resolución)
\begin{frame}{Ejercicio en Clase --- Nivel 2: Operativo (2/2)}
  \begin{columns}[T]
    \begin{column}{0.48\textwidth}
      \begin{block}{1. Valor Predictivo Positivo (VPP)}
        \fontsize{6.2pt}{7.4pt}\selectfont
        Calculamos la probabilidad marginal total de un resultado positivo $P(+)$:
        $$P(+) = P(+ \mid E)P(E) + P(+ \mid E')P(E')$$
        $$P(+) = (0.95)(0.02) + (0.10)(0.98) = 0.019 + 0.098 = 0.117$$
        Aplicando Bayes para hallar el VPP ($P(E \mid +)$):
        $$P(E \mid +) = \frac{0.019}{0.117} \approx \mathbf{0.16239 \text{ (16.24\%)}}$$
      \end{block}
    \end{column}
    \begin{column}{0.48\textwidth}
      \begin{alertblock}{2. Valor Predictivo Negativo (VPN)}
        \fontsize{6.2pt}{7.4pt}\selectfont
        Calculamos la probabilidad marginal de un resultado negativo $P(-) = 1 - P(+) = 0.883$:
        $$P(E' \mid -) = \frac{P(- \mid E')P(E')}{P(-)}$$
        $$P(E' \mid -) = \frac{(0.90)(0.98)}{0.883} = \frac{0.882}{0.883} \approx \mathbf{0.99887 \text{ (99.89\%)}}$$
      \end{alertblock}
    \end{column}
  \end{columns}
  \vspace{-0.1cm}
  \pause
  \begin{exampleblock}{Análisis de la Falacia de la Tasa Base}
    \fontsize{6.2pt}{7.4pt}\selectfont
    ¡El $83.76\%$ de los pacientes que obtienen un resultado positivo son falsas alarmas! Esto no se debe a que el test sea inútil ($95\%$ de sensibilidad), sino a que la masa demográfica de personas sanas ($98\%$) multiplicada por su tasa de error ($10\%$) genera un volumen de falsos positivos ($0.098$) que supera 5 a 1 a los verdaderos positivos ($0.019$).
  \end{exampleblock}
\end{frame}

% Slide 12A: Ejercicio en Clase — Nivel 3: Analítico (Enunciado)
\begin{frame}{Ejercicio en Clase --- Nivel 3: Analítico (1/2)}
  \begin{block}{Problema 2.5.3 --- Auditoría Bayesiana en Líneas de Ensamble}
    \footnotesize
    Una planta de semiconductores opera con tres líneas de manufactura simultáneas ($M_1, M_2, M_3$) que aportan el $50\%$, $30\%$ y $20\%$ de la producción total, respectivamente. Las tasas de defectos estructurales ($D$) verificadas en control de calidad son:
    $$P(D \mid M_1) = 0.03, \qquad P(D \mid M_2) = 0.05, \qquad P(D \mid M_3) = 0.07$$
    \begin{enumerate}\itemsep=0.15cm
      \item Si un auditor externo selecciona un chip defectuoso de los contenedores finales, determina el vector de probabilidades a posteriori $P(M_i \mid D)$ para identificar el origen de la falla.
      \item Demuestra analíticamente cuál línea es estadísticamente la causa más probable del defecto y explica el equilibrio entre volumen de producción y tasa de error.
    \end{enumerate}
  \end{block}
\end{frame}

% Slide 12B: Ejercicio en Clase — Nivel 3: Analítico (Resolución)
\begin{frame}{Ejercicio en Clase --- Nivel 3: Analítico (2/2)}
  \begin{columns}[T]
    \begin{column}{0.48\textwidth}
      \begin{block}{1. Tasa Marginal y Vector a Posteriori}
        \fontsize{6.2pt}{7.4pt}\selectfont
        Calculamos la probabilidad total de defecto en la planta:
        $$P(D) = \sum_{i=1}^3 P(D \mid M_i)P(M_i)$$
        $$P(D) = (0.03)(0.50) + (0.05)(0.30) + (0.07)(0.20)$$
        $$P(D) = 0.015 + 0.015 + 0.014 = \mathbf{0.044 \text{ (4.4\%)}}$$
        Invertimos con Bayes para cada línea $M_i$:
        \begin{itemize}\itemsep=0.04cm
          \item $P(M_1 \mid D) = \frac{0.015}{0.044} = \frac{15}{44} \approx \mathbf{34.09\%}$.
          \item $P(M_2 \mid D) = \frac{0.015}{0.044} = \frac{15}{44} \approx \mathbf{34.09\%}$.
          \item $P(M_3 \mid D) = \frac{0.014}{0.044} = \frac{14}{44} \approx \mathbf{31.82\%}$.
        \end{itemize}
      \end{block}
    \end{column}
    \begin{column}{0.48\textwidth}
      \begin{alertblock}{2. Análisis de Empate y Recomendación de Auditoría}
        \fontsize{6.2pt}{7.4pt}\selectfont
        \textbf{Conclusión:} Tanto la línea $M_1$ como la línea $M_2$ empatan exactamente con un $34.09\%$ de probabilidad como la causa más probable de cualquier componente defectuoso detectado en el inventario final.
        \vspace{0.1cm}
        
        \textbf{Justificación Algorítmica:} Aunque la línea $M_1$ posee la tasa de defectos más baja ($3\%$) en comparación con $M_2$ ($5\%$), su volumen masivo de manufactura ($50\%$ vs $30\%$) genera exactamente la misma cantidad absoluta de piezas fallidas en el flujo total:
        $$P(D \cap M_1) = P(D \cap M_2) = 0.015$$
      \end{alertblock}
    \end{column}
  \end{columns}
\end{frame}

% Slide 13A: Ejercicio en Clase — Nivel 4: Desafiante (Enunciado)
\begin{frame}{Ejercicio en Clase --- Nivel 4: Desafiante (1/2)}
  \begin{block}{Problemas 2.5.9 \& 2.5.10 --- Falacia del Fiscal y Filtro Naive Bayes}
    \footnotesize
    \begin{itemize}\itemsep=0.15cm
      \item \textbf{Falacia del Fiscal (2.5.9):} En un caso forense, una muestra biológica coincide con el perfil de ADN de un sospechoso ($E$). La tasa de coincidencia aleatoria en una población de $N=5,000,000$ es de $\alpha=10^{-6}$. El fiscal afirma ante el juez que el sospechoso tiene $99.9999\%$ de probabilidad de culpabilidad. Utiliza la razón de momios bayesiana para calcular la verdadera probabilidad a posteriori $P(\text{Culpable} \mid E)$.
      \item \textbf{Clasificador Naive Bayes Anti-Spam (2.5.10):} Un servidor procesa $60\%$ de correos spam ($S$) y $40\%$ legítimos ($S'$). Bajo independencia condicional, evalúa un mensaje que contiene las palabras clave $W_1 \cap W_2 \cap W_3$ con verosimilitudes: $P(W \mid S)=(0.80)(0.70)(0.60)=0.336$ y $P(W \mid S')=(0.10)(0.20)(0.10)=0.002$. Halla $P(S \mid W)$.
    \end{itemize}
  \end{block}
\end{frame}

% Slide 13B: Ejercicio en Clase — Nivel 4: Desafiante (Resolución)
\begin{frame}{Ejercicio en Clase --- Nivel 4: Desafiante (2/2)}
  \vspace{-0.15cm}
  \begin{columns}[T]
    \begin{column}{0.48\textwidth}
      \begin{block}{1. Falacia del Fiscal en Momios ($N=5 \times 10^6$)}
        \fontsize{5.5pt}{6.4pt}\selectfont
        Momios a priori en una población de 5 millones de sospechosos potenciales:
        $$\text{Momios}(C) = \frac{1 / 5,000,000}{1} = 2 \times 10^{-7}$$
        Cociente de Verosimilitud ($\text{LR}$) ante coincidencia de ADN:
        $$\text{LR} = \frac{P(E \mid C)}{P(E \mid C')} = \frac{1}{10^{-6}} = 1,000,000$$
        Momios a posteriori y probabilidad real:
        $$\text{Momios}(C \mid E) = 10^6 \times (2 \times 10^{-7}) = 0.20 = \frac{1}{5}$$
        $$P(C \mid E) = \frac{0.20}{1 + 0.20} = \frac{1}{6} \approx \mathbf{16.67\%}$$
        ¡La afirmación del fiscal ($99.9999\%$) es una falacia por ignorar la población marginal!
      \end{block}
    \end{column}
    \begin{column}{0.48\textwidth}
      \begin{alertblock}{2. Inversión Naive Bayes en Spam ($W$)}
        \fontsize{5.5pt}{6.4pt}\selectfont
        Calculamos la probabilidad marginal conjunta del vector de palabras $W$:
        $$P(W) = P(W \mid S)P(S) + P(W \mid S')P(S')$$
        $$P(W) = (0.336)(0.60) + (0.002)(0.40) = 0.2016 + 0.0008 = 0.2024$$
        Calculamos la probabilidad a posteriori exacta de que el mensaje sea spam ($S$):
        $$P(S \mid W) = \frac{0.2016}{0.2024} = \frac{252}{253} \approx \mathbf{99.6047\%}$$
        La combinación simultánea de los tres indicadores clasifica el correo con cuasi-certeza.
      \end{alertblock}
    \end{column}
  \end{columns}
\end{frame}

% Slide 14: Síntesis y Conclusión
\begin{frame}{Síntesis y Conclusión --- Teorema de Bayes e Inversión}
  \fontsize{7.5pt}{9pt}\selectfont
  \begin{block}{Pilares Conceptuales de la Sección 02.05}
    \fontsize{7.5pt}{9pt}\selectfont
    \begin{itemize}\itemsep=0.12cm
      \item \textbf{Inversión Rigurosa de Condicionales:} El Teorema de Bayes proporciona el mecanismo algebraico exacto para inferir probabilidades de causas ($P(A \mid B)$) a partir de verosimilitudes empíricas ($P(B \mid A)$).
      \item \textbf{El Peso Crítico del Prior:} Ignorar la probabilidad marginal a priori $P(A)$ conduce inevitablemente a sesgos cognitivos graves y a la falacia de la tasa base en sistemas de diagnóstico técnico y clínico.
      \item \textbf{Actualización Iterativa y Robótica:} La transición donde la distribución a posteriori de un ciclo se convierte en la creencia a priori del siguiente ciclo es el fundamento de los algoritmos de IA y estimación en tiempo real.
    \end{itemize}
  \end{block}
\end{frame}

% Slide 15: Transición Curricular
\begin{frame}{Transición Curricular y Conexión Temática}
  \fontsize{7.5pt}{9pt}\selectfont
  \begin{block}{Conexión Curricular Hacia Atrás y Hacia Adelante}
    \fontsize{7.5pt}{9pt}\selectfont
    Con el dominio del Teorema de Bayes, hemos completado las herramientas analíticas para manipular probabilidades sobre eventos conocidos o fijos en el espacio muestral.
    \begin{itemize}\itemsep=0.1cm
      \item La probabilidad condicional y la inversión de Bayes nos permiten evaluar hipótesis específicas dadas pruebas observacionales concretas.
      \item Sin embargo, para generalizar esta inferencia a muestras extraídas aleatoriamente de grandes poblaciones, necesitamos estudiar el comportamiento sintético y asintótico de los promedios muestrales.
    \end{itemize}
  \end{block}
  \vspace{-0.15cm}
  \pause
  \begin{alertblock}{Conexión con la Sección 02.06: Muestreo Aleatorio y TLC}
    \fontsize{7.5pt}{9pt}\selectfont
    En la \textbf{Sección 02.06: Muestreo Aleatorio y Teorema del Límite Central}, exploramos cómo las distribuciones muestrales de variables aleatorias independientes convergen algebraicamente hacia la normalidad universal, permitiendo acotar el error estadístico con precisión garantizada.
  \end{alertblock}
\end{frame}

\end{document}
